\subsection{A cache on one affine line: every temperature}
\label{sec:newrankone}

A cache on one affine line permits exact sampling whose cost
depends only on the encoding length of the temperature, even when
the numerical temperature is very large.
The proof uses the geometric rejection envelope for decreasing
weights studied by Chewi, Gerber, Lu, Le Gouic, and
Rigollet~\cite{chewi2022}. We include the short envelope argument
and account separately for finite arithmetic.

\begin{lemma}[A geometric partition of decreasing weights]
\label{lem:newmonotoneenvelope}
Let $w_1\ge\cdots\ge w_T>0$ and $Z=\sum_jw_j$.
Partition $[T]$ into $G_0=\{1\}$ and the nonempty blocks
\[
 G_\ell=\{2^{\ell-1}+1,\ldots,\min(2^\ell,T)\},\qquad\ell\ge1.
\]
Write $m_\ell=\min G_\ell$, $N_\ell=|G_\ell|$, and
$W=\sum_\ell N_\ell w_{m_\ell}$. Then
\[
 Z\le W\le2Z.
\]
Thus sampling a block proportionally to $N_\ell w_{m_\ell}$,
sampling a uniform member $j$, and accepting it with probability
$w_j/w_{m_\ell}$ gives an exact sample from $w_j/Z$ with at most
two proposals in expectation.
\end{lemma}
\paragraph{Proof idea.}
Each geometric block borrows its envelope mass from the preceding block.  Every true weight is charged at most twice, giving a constant proposal count for any decreasing sequence.

\begin{proof}
The lower bound follows from monotonicity within each block. If
$T\le2$ the envelope is exact. Otherwise, for every $\ell\ge2$
with a nonempty block, its preceding block is full, has size
$2^{\ell-2}$, and all its weights are at least $w_{m_\ell}$.
Consequently
\[
 N_\ell w_{m_\ell}
 \le2\sum_{j\in G_{\ell-1}}w_j.
\]
If $L$ is the last block, summation gives
\[
 W\le w_1+w_2+2\sum_{\ell=1}^{L-1}\sum_{j\in G_\ell}w_j
 =2Z-w_1+w_2-2\sum_{j\in G_L}w_j\le2Z.
\]
The accepted mass in a trial is $w_j/W$, and the success
probability is $Z/W$. Summing the geometric series proves both
exactness and the expected proposal count.
\end{proof}

\begin{lemma}[Exponential evaluation on the negative half-line]
\label{lem:newnegativeexp}
Suppose $x\le0$ is either a rational with encoding length $b$, or
a rational plus a rational divided by $\sqrt n$, with total
encoding length $b$.
For every integer $p\ge1$, a deterministic finite algorithm returns
a dyadic interval containing $e^x$ of length at most $2^{-p}$ in
$O((b+p+\log(n+2))^4)$ bit operations. The bound has no dependence
on the numerical magnitude $|x|$ beyond its encoding length.
\end{lemma}
\paragraph{Proof idea.}
An exponential that is already smaller than the requested error can be enclosed immediately.  Every remaining argument lies in a precision-dependent bounded interval, where a finite Taylor calculation and outward rounding suffice.

\begin{proof}
First obtain an interval of length at most $1/4$ for $x$.
Rational arithmetic, and rational bisection for $\sqrt n$ when
needed, have the asserted polynomial cost. If the upper endpoint
is at most $-(p+2)$, return $[0,2^{-(p+2)}]$; this is valid because
$e>2$. Otherwise the upper endpoint is greater than $-(p+2)$ and
the interval has width at most $1/4$, so the known sign gives
$|x|\le p+3$. Refine $x$ to absolute error
$2^{-(p+3)}$ and evaluate its exponential on $[-p-3,0]$.
One explicit procedure divides the argument by
$m=p+4$, uses the alternating exponential series on $[-1,0]$,
and raises the resulting enclosure to the power $m$.
Take $K=4(p+\lceil\log_2(m+1)\rceil+10)$ terms. The alternating
remainder is at most $1/(K+1)!\le2^{-K}$.
Perform interval arithmetic with outward rounding to
$P=p+2\lceil\log_2(K+1)\rceil+\lceil\log_2(m+1)\rceil+20$
fractional bits. In the term recurrence, multiplication by a
number of absolute value at most one and division by a positive
integer do not increase the preceding absolute error. Summing
the $K$ terms therefore adds at most $4(K+1)^2 2^{-P}$ of
rounding error. Clip the resulting exponential enclosure to
$[0,1]$ and multiply it by itself $m$ times, again clipping to
$[0,1]$. These products increase interval width by at most the
sum of their input widths plus $2^{1-P}$ per product.
The chosen $K,P$ make the final error at most $2^{-(p+2)}$.
Monotonicity and $|\exp'(x)|\le1$ on the negative half-line
control the input enclosure error. There are $O(p+\log(p+2))$
fixed-point operations on $O(b+p+\log n)$ bits; schoolbook
arithmetic is bounded by the fourth power displayed above.
In the algebraic case, an absolute interval
for the rational numerator divided by $\sqrt n$ needs only a
linear number of extra precision bits in the numerator's encoding
length, even if the numerator has large magnitude.
\end{proof}

\begin{theorem}[Exact indexed attention in affine dimension one]
\label{thm:newrankone}
Let $T\ge2$. Let cached keys $k_j\in\mathbb Q^n$ lie on one
affine line, with arbitrary cached scalar values $v_j\in[-1,1]$.
All keys, values, and the fresh query $q\in\mathbb Q^n$ have dyadic
coordinates with at most $b$ bits per coordinate. Let $\beta\ge0$
be dyadic, also with at most $b$ bits, or let $\beta=\sqrt n$.
There is a uniform algorithm with the following bounds:
\[
 \begin{split}
 \text{prefill work}&=O\bigl((Tn+T\log T)(b+\log(T+n))^2\bigr),\\
 \text{stored bits}&=O\bigl(Tn b+T\log T\bigr),\\
 \text{expected online work}&=
 O\bigl((n+\log T)(b+\log(T+n))^4\bigr).
 \end{split}
\]
It samples an exact attention index for scores
$\beta\langle q,k_j\rangle/n$, and hence an exact sign of mean
$\sum_j e^{a_j}v_j/\sum_j e^{a_j}$.
Its expected number of proposed cache members is at most $17/8$.
These bounds hold at every numerical temperature covered by the
stated finite encoding. No upper bound on the scores is required.
\end{theorem}
\paragraph{Proof idea.}
Project the query onto the cache line once.  The projection chooses one of two stored orders; a geometric envelope then samples from that order.  A small dyadic floor makes every acceptance comparison efficient even at very large temperature.

\begin{proof}
If all keys coincide, attention is uniform and the claim follows
immediately. Otherwise choose two distinct keys, set $c=k_1$, and
let $u$ be their nonzero difference. Choose a pivot coordinate $r$
with $u_r\ne0$. Prefill can verify the affine-line promise by
checking
\[
 u_r(k_{j,i}-c_i)=u_i(k_{j,r}-c_r)
 \quad\text{for every }i,j.
\]
Sort the keys by their pivot coordinate, and store that permutation.
Also store the rational affine coordinates
$\xi_j=(k_{j,r}-c_r)/u_r$, so $k_j=c+\xi_ju$.
All these tests and comparisons use exact rational arithmetic.

At query time, compute $v=\langle q,u\rangle$ once and the sign of $v/u_r$.
It determines the order of all projected keys; use the stored
permutation in the appropriate direction. If the sign is zero,
all scores coincide. Otherwise relabel the keys in descending
score order and shift by the first score. Set
\[
 w_j=\exp(a_j-a_1),\qquad 1=w_1\ge\cdots\ge w_T>0.
\]
Each shift is computed as $\beta v(\xi_j-\xi_1)/n$.
It is rational, or a rational multiple of $1/\sqrt n$ at standard
scaling. Its encoding has $O(b+\log n)$ bits. Thus one initial
projection suffices for every subsequent score evaluation, and no
subtraction of numerically approximated scores is required.
Use the blocks of Lemma~\ref{lem:newmonotoneenvelope}; their number
is at most $1+\lceil\log_2T\rceil$.

Put $P=\lceil\log_2(16T)\rceil$. For each first member of a block,
Lemma~\ref{lem:newnegativeexp} gives a dyadic interval of length at
most $2^{-P}$ for $w_{m_\ell}$. Add $2^{-P}$ to its upper endpoint,
and call the result $u_\ell$. Then
\[
 2^{-P}\le u_\ell,
 \qquad w_{m_\ell}\le u_\ell\le w_{m_\ell}+2^{1-P}.
\]
The actual dyadic envelope satisfies
\[
 Z\le W':=\sum_\ell N_\ell u_\ell
 \le2Z+T2^{1-P}\le2Z+\tfrac18\le\tfrac{17}{8}Z,
\]
where $Z\ge w_1=1$. Its weights have $O(P)$ bits after a common
dyadic scaling, and exact integer prefix sums sample a block in
$O((P+\log T)^2)$ expected bit operations. Select a uniform member
$j$ of that block and accept it with probability $w_j/u_\ell$.
This lies in $[0,1]$, and the trial's accepted mass is $w_j/W'$.
Thus the index law is exact and the mean number of trials is
$W'/Z\le17/8$.

To implement the last coin to precision $t$, it suffices to enclose
$w_j$ to precision $P+t+3$, since $u_\ell\ge2^{-P}$.
The negative-exponential lemma does this in
$O((b+P+t+\log(n+2))^4)$ bit work, regardless of the temperature's
numerical size. Lemma~\ref{lem:newprecisionstop} therefore makes the
coin exact with expected work $O((b+P+\log(n+2))^4)$. There is one
$n$-coordinate projection, followed by $O(\log T)$ representative
score evaluations and a constant expected number of proposed-member
score evaluations. A fresh known
coin of mean $v_j$ completes the attention sign.

The preprocessing bounds follow from the collinearity tests and
sorting. Rational bit lengths remain $O(b+\log n)$ up to constant
factors for each comparison and inner product. The stored
permutation, block positions, and exact dyadic coordinates have
the stated total size.
\end{proof}

\paragraph{What is being sampled.}
The theorem gives a sample from the attention distribution and a
sign with the exact attention-coordinate mean. A nonlinearity applied
to one sampled value vector does not in general give the required
law. An additional nonlinear layer instead requests fresh
independent attention signs through its Bernoulli factory. For
example, a tanh update followed by a two-logit head $(z,-z)$ with
$z=\tanh Y$ needs only a constant expected number of these signs,
using the bounded scalar backend. Thus the same asymptotic query
bound applies to that exact final token.
